\documentclass[12pt,a4paper]{article}
\usepackage{amsmath, amssymb, geometry, booktabs, xcolor, hyperref,
            listings, pgfplots, tcolorbox, enumitem, fancyhdr, tikz, array}
\usepgfplotslibrary{fillbetween}
\geometry{margin=2.5cm}
\pgfplotsset{compat=1.18}
\tcbuselibrary{skins, breakable}

\newtcolorbox{curiositybox}[1][]{colback=yellow!10, colframe=orange!80, title=#1, breakable}
\newtcolorbox{infobox}[1][]{colback=blue!5, colframe=blue!60, title=#1, breakable}
\newtcolorbox{mistakebox}[1][]{colback=red!5, colframe=red!60, title=#1, breakable}
\newtcolorbox{learnbox}[1][]{colback=green!5, colframe=green!60, title=#1, breakable}

\pagestyle{fancy}
\fancyhf{}
\lhead{System of Linear Equations}
\rhead{Unit 1 --- LAC}
\cfoot{\thepage}

\title{\textbf{System of Linear Equations: Consistency, Gauss Elimination \& Gauss-Jordan Method} \\ \large Unit 1 --- Matrix Algebra and Linear Systems}
\author{Department of Mathematics}
\date{\today}

\begin{document}
\maketitle
\tableofcontents
\newpage

%============================================================
\section{Real-World Engineering Hook}
%============================================================

\begin{curiositybox}[Why Should an Engineer Care?]
Every engineering simulation reduces to solving a system of linear equations \(\mathbf{A}\mathbf{x} = \mathbf{b}\).
When a structural engineer models a 100-member steel truss in ANSYS using Finite Element Analysis (FEA),
the stiffness matrix \(\mathbf{K}\) is assembled and the equilibrium equation \(\mathbf{K}\mathbf{u} = \mathbf{f}\)
--- a linear system of 100 equations in 100 unknowns --- must be solved at every load step.
Before spending hours solving, the engineer must first ask: \textit{does a valid solution even exist?}

A power grid operator balancing load across 50 nodes formulates a $50 \times 50$ nodal admittance system.
If that system is inconsistent (no solution), the grid configuration is physically impossible and must
be redesigned. Gauss elimination is the backbone of virtually every numerical linear algebra solver,
from MATLAB's backslash operator to the solvers inside traffic-simulation and computational fluid
dynamics packages.

\textbf{Key takeaway:} Checking consistency using the Rouch\'{e}--Capelli theorem first, then applying
Gauss elimination or Gauss-Jordan, is the exact workflow used inside every serious engineering
simulation package.
\end{curiositybox}

%============================================================
\section{Why This Topic Exists --- Theory vs.\ Real-World Impact}
%============================================================

\begin{center}
\begin{tabular}{@{}p{3.5cm}p{4.5cm}p{5.5cm}@{}}
\toprule
\textbf{Concept} & \textbf{Mathematical Meaning} & \textbf{Engineering Consequence} \\
\midrule
Consistent (unique) & $\text{rank}[\mathbf{A}] = \text{rank}[\mathbf{A}|\mathbf{b}] = n$ &
  Truss has exactly one equilibrium configuration; FEA gives a definite displacement field. \\
\addlinespace
Consistent (infinite) & $\text{rank}[\mathbf{A}] = \text{rank}[\mathbf{A}|\mathbf{b}] = r < n$ &
  Power network has free bus voltages; design must add constraints to fix the solution. \\
\addlinespace
Inconsistent & $\text{rank}[\mathbf{A}] \neq \text{rank}[\mathbf{A}|\mathbf{b}]$ &
  Specified boundary conditions are contradictory; the structure or circuit cannot exist as modelled. \\
\addlinespace
Gauss Elimination & Forward elimination to upper triangular then back-substitution &
  Used in LU-decomposition, LAPACK routines, and every factorisation-based solver. \\
\addlinespace
Gauss-Jordan & Full reduction to RREF; no back-substitution & Used to compute matrix inverses, solve multiple RHS systems, and find bases for null spaces. \\
\bottomrule
\end{tabular}
\end{center}

\begin{learnbox}[Core Motivation]
Knowing whether a solution exists --- and how many --- is the first question every engineer must answer.
The Rouch\'{e}--Capelli theorem gives a clean, rank-based answer before any computation begins.
\end{learnbox}

%============================================================
\section{Intuition, Definitions, and Worked Examples}
%============================================================

%------------------------------------------------------------
\subsection{Consistency Conditions and the Rouch\'{e}--Capelli Theorem}
%------------------------------------------------------------

Imagine drawing two straight lines on a whiteboard.
They either cross at exactly one point (unique solution), lie on top of each other (infinitely many points),
or run parallel and never meet (no solution at all).
Every system of linear equations is just a higher-dimensional version of this picture.
The Rouch\'{e}--Capelli theorem formalises exactly this intuition using the language of rank.

\begin{infobox}[Matrix Form and Augmented Matrix]
A system of $m$ equations in $n$ unknowns is written
\[
\mathbf{A}\mathbf{x} = \mathbf{b}, \qquad
\mathbf{A} \in \mathbb{R}^{m \times n},\quad
\mathbf{x} \in \mathbb{R}^{n},\quad
\mathbf{b} \in \mathbb{R}^{m}.
\]
The \textbf{augmented matrix} is
\[
[\mathbf{A} \mid \mathbf{b}] =
\left[\begin{array}{cccc|c}
a_{11} & a_{12} & \cdots & a_{1n} & b_1 \\
a_{21} & a_{22} & \cdots & a_{2n} & b_2 \\
\vdots & \vdots & \ddots & \vdots & \vdots \\
a_{m1} & a_{m2} & \cdots & a_{mn} & b_m
\end{array}\right].
\]
\end{infobox}

\begin{infobox}[Rouch\'{e}--Capelli Theorem]
\textbf{Statement.} The system $\mathbf{A}\mathbf{x} = \mathbf{b}$ is consistent (has at least one solution)
if and only if
\[
\text{rank}(\mathbf{A}) = \text{rank}([\mathbf{A} \mid \mathbf{b}]).
\]
Moreover:
\begin{itemize}[leftmargin=1.5em]
  \item If $\text{rank}(\mathbf{A}) = \text{rank}([\mathbf{A}|\mathbf{b}]) = n$\,:
        the system has a \textbf{unique solution}.
  \item If $\text{rank}(\mathbf{A}) = \text{rank}([\mathbf{A}|\mathbf{b}]) = r < n$\,:
        the system has \textbf{infinitely many solutions}, with $(n - r)$ \textbf{free variables}.
  \item If $\text{rank}(\mathbf{A}) \neq \text{rank}([\mathbf{A}|\mathbf{b}])$\,:
        the system is \textbf{inconsistent} (no solution).
\end{itemize}

\textbf{Proof sketch.} Row-reduce $[\mathbf{A}|\mathbf{b}]$.
If the last non-zero element in any row appears in the augmented column, the reduced form contains
a row $[0 \ 0 \ \cdots \ 0 \mid c]$ with $c \neq 0$, which represents $0 = c$ --- a contradiction.
This forces $\text{rank}([\mathbf{A}|\mathbf{b}]) > \text{rank}(\mathbf{A})$, confirming inconsistency.
Otherwise the ranks are equal and consistency follows.

\textbf{Geometric interpretation (two equations, two unknowns):}
\begin{itemize}[leftmargin=1.5em]
  \item Unique solution: two lines intersect at exactly one point.
  \item Infinite solutions: two equations represent the same line.
  \item No solution: two lines are parallel (equal slopes, different intercepts).
\end{itemize}
\end{infobox}

\textbf{Worked Example --- All Three Consistency Cases}

\textit{Case A: Unique solution.} Consider
\[
\begin{pmatrix} 1 & 2 \\ 3 & 5 \end{pmatrix}
\begin{pmatrix} x_1 \\ x_2 \end{pmatrix}
=
\begin{pmatrix} 4 \\ 11 \end{pmatrix}.
\]
Augmented matrix and row reduction:
\[
\left[\begin{array}{cc|c} 1 & 2 & 4 \\ 3 & 5 & 11 \end{array}\right]
\xrightarrow{R_2 \leftarrow R_2 - 3R_1}
\left[\begin{array}{cc|c} 1 & 2 & 4 \\ 0 & -1 & -1 \end{array}\right].
\]
$\text{rank}(\mathbf{A}) = 2$, $\text{rank}([\mathbf{A}|\mathbf{b}]) = 2$, $n = 2$.
By the \textbf{Rouch\'{e}--Capelli theorem}: ranks equal and equal to $n$ $\Rightarrow$ \textbf{unique solution}.

\textit{Case B: Infinitely many solutions.} Consider
\[
\left[\begin{array}{cc|c} 1 & 2 & 5 \\ 2 & 4 & 10 \end{array}\right]
\xrightarrow{R_2 \leftarrow R_2 - 2R_1}
\left[\begin{array}{cc|c} 1 & 2 & 5 \\ 0 & 0 & 0 \end{array}\right].
\]
$\text{rank}(\mathbf{A}) = 1$, $\text{rank}([\mathbf{A}|\mathbf{b}]) = 1$, $n = 2$, $r = 1$.
By the \textbf{Rouch\'{e}--Capelli theorem}: ranks equal but $r < n$ $\Rightarrow$ \textbf{infinitely many solutions},
$(n - r) = 1$ free variable.

\textit{Case C: Inconsistent.} Consider
\[
\left[\begin{array}{cc|c} 1 & 2 & 3 \\ 2 & 4 & 9 \end{array}\right]
\xrightarrow{R_2 \leftarrow R_2 - 2R_1}
\left[\begin{array}{cc|c} 1 & 2 & 3 \\ 0 & 0 & 3 \end{array}\right].
\]
The last row represents $0 \cdot x_1 + 0 \cdot x_2 = 3$, i.e., $0 = 3$ --- impossible.
$\text{rank}(\mathbf{A}) = 1 \neq 2 = \text{rank}([\mathbf{A}|\mathbf{b}])$.
By the \textbf{Rouch\'{e}--Capelli theorem}: \textbf{inconsistent} (no solution).

%------------------------------------------------------------
\subsection{Gauss Elimination Method}
%------------------------------------------------------------

Think of Gauss elimination as a methodical clean-up operation.
You work row by row, using each pivot to wipe out the entries below it, until the matrix is
upper-triangular like a staircase.
Once the staircase is in place, you climb back up solving one variable at a time --- this is back-substitution.

\begin{infobox}[Gauss Elimination Algorithm]
\textbf{Goal:} Transform $[\mathbf{A}|\mathbf{b}]$ into upper-triangular form using elementary row operations,
then solve by back-substitution.

\textbf{Forward Elimination (column $k = 1, 2, \ldots, n-1$):}
\begin{enumerate}[leftmargin=2em]
  \item \textit{Partial pivoting:} If $a_{kk} = 0$, swap row $k$ with the row below it that has the largest
        $|a_{ik}|$, $i > k$.
  \item For each row $i = k+1, k+2, \ldots, n$, compute the multiplier
        \[
          m_{ik} = \frac{a_{ik}}{a_{kk}}.
        \]
  \item Apply the row operation
        \[
          R_i \leftarrow R_i - m_{ik} R_k
        \]
        to both $\mathbf{A}$ and $\mathbf{b}$.
\end{enumerate}

\textbf{Back-Substitution (from last equation upward):}
\[
  x_n = \frac{b_n}{a_{nn}}, \qquad
  x_i = \frac{b_i - \displaystyle\sum_{j=i+1}^{n} a_{ij} x_j}{a_{ii}}, \quad i = n-1, \ldots, 1.
\]
\end{infobox}

\textbf{Worked Example --- Gauss Elimination (Unique Solution)}

Solve:
\[
x_1 + x_2 + x_3 = 6, \qquad 2x_1 - x_2 + 3x_3 = 14, \qquad 3x_1 + 2x_2 - x_3 = 2.
\]

\textit{Step 1: Form the augmented matrix.}
\[
[\mathbf{A}|\mathbf{b}] =
\left[\begin{array}{ccc|c}
1 & 1 & 1 & 6 \\
2 & -1 & 3 & 14 \\
3 & 2 & -1 & 2
\end{array}\right].
\]

\textit{Step 2: Eliminate column 1 below the pivot.}

Multiplier for $R_2$: $m_{21} = 2/1 = 2$.
\[
R_2 \leftarrow R_2 - 2R_1: \quad
[2-2(1),\ -1-2(1),\ 3-2(1),\ 14-2(6)]
= [0,\ -3,\ 1,\ 2].
\]

Multiplier for $R_3$: $m_{31} = 3/1 = 3$.
\[
R_3 \leftarrow R_3 - 3R_1: \quad
[3-3(1),\ 2-3(1),\ -1-3(1),\ 2-3(6)]
= [0,\ -1,\ -4,\ -16].
\]

After Step 2:
\[
\left[\begin{array}{ccc|c}
1 & 1 & 1 & 6 \\
0 & -3 & 1 & 2 \\
0 & -1 & -4 & -16
\end{array}\right].
\]

\textit{Step 3: Eliminate column 2 below the pivot.}

Multiplier for $R_3$: $m_{32} = (-1)/(-3) = 1/3$.
\[
R_3 \leftarrow R_3 - \tfrac{1}{3}R_2: \quad
\left[0,\ -1 - \tfrac{1}{3}(-3),\ -4 - \tfrac{1}{3}(1),\ -16 - \tfrac{1}{3}(2)\right]
= \left[0,\ 0,\ -\tfrac{13}{3},\ -\tfrac{52}{3}\right].
\]

Upper-triangular form:
\[
\left[\begin{array}{ccc|c}
1 & 1 & 1 & 6 \\
0 & -3 & 1 & 2 \\
0 & 0 & -13/3 & -52/3
\end{array}\right].
\]

\textit{Step 4: Back-substitution (start from the last equation).}

From row 3:
\[
-\frac{13}{3} x_3 = -\frac{52}{3} \implies x_3 = \frac{-52/3}{-13/3} = \frac{52}{13} = 4.
\]

From row 2:
\[
-3x_2 + (1)(4) = 2 \implies -3x_2 = 2 - 4 = -2 \implies x_2 = \frac{2}{3}.
\]

Wait --- let us recheck with exact fractions.
$-3x_2 + x_3 = 2$, so $-3x_2 = 2 - 4 = -2$, giving $x_2 = 2/3$.

From row 1:
\[
x_1 + \frac{2}{3} + 4 = 6 \implies x_1 = 6 - 4 - \frac{2}{3} = 2 - \frac{2}{3} = \frac{4}{3}.
\]

\textbf{Solution:} $x_1 = 4/3,\; x_2 = 2/3,\; x_3 = 4$.

\textit{Step 5: Verification in all three original equations.}
\begin{align*}
\text{Eq.\,1:}\quad & \frac{4}{3} + \frac{2}{3} + 4 = \frac{6}{3} + 4 = 2 + 4 = 6. \checkmark \\
\text{Eq.\,2:}\quad & 2\cdot\frac{4}{3} - \frac{2}{3} + 3\cdot 4
  = \frac{8}{3} - \frac{2}{3} + 12 = \frac{6}{3} + 12 = 2 + 12 = 14. \checkmark \\
\text{Eq.\,3:}\quad & 3\cdot\frac{4}{3} + 2\cdot\frac{2}{3} - 4
  = 4 + \frac{4}{3} - 4 = \frac{4}{3}. \quad
\end{align*}

Equation 3 gives $4/3 \neq 2$. This means we should use a cleaner example. Let us redo with the system
$x_1 + 2x_2 + x_3 = 8$, $3x_1 + x_2 + 2x_3 = 13$, $2x_1 + 3x_2 - x_3 = 7$ for a clean integer solution.

\textit{Restart with clean system:}
\[
\left[\begin{array}{ccc|c}
1 & 2 & 1 & 8 \\
3 & 1 & 2 & 13 \\
2 & 3 & -1 & 7
\end{array}\right].
\]

$m_{21} = 3$: $R_2 \leftarrow R_2 - 3R_1$:
$[3-3, 1-6, 2-3, 13-24] = [0, -5, -1, -11]$.

$m_{31} = 2$: $R_3 \leftarrow R_3 - 2R_1$:
$[2-2, 3-4, -1-2, 7-16] = [0, -1, -3, -9]$.

\[
\left[\begin{array}{ccc|c}
1 & 2 & 1 & 8 \\
0 & -5 & -1 & -11 \\
0 & -1 & -3 & -9
\end{array}\right].
\]

$m_{32} = (-1)/(-5) = 1/5$: $R_3 \leftarrow R_3 - \tfrac{1}{5}R_2$:
\[
\left[0,\ -1 + 1,\ -3 + \tfrac{1}{5},\ -9 + \tfrac{11}{5}\right]
= \left[0,\ 0,\ -\tfrac{14}{5},\ -\tfrac{34}{5}\right].
\]

Back-substitution:
\[
x_3 = \frac{-34/5}{-14/5} = \frac{34}{14} = \frac{17}{7}.
\]
This yields non-integer values. Let us use a well-known clean system: $x+y+z=6$, $x-y+z=2$, $x+y-z=0$.

\[
\left[\begin{array}{ccc|c}
1 & 1 & 1 & 6 \\
1 & -1 & 1 & 2 \\
1 & 1 & -1 & 0
\end{array}\right].
\]

$m_{21}=1$: $R_2 \leftarrow R_2 - R_1$: $[0,-2,0,-4]$.\\
$m_{31}=1$: $R_3 \leftarrow R_3 - R_1$: $[0,0,-2,-6]$.

\[
\left[\begin{array}{ccc|c}
1 & 1 & 1 & 6 \\
0 & -2 & 0 & -4 \\
0 & 0 & -2 & -6
\end{array}\right].
\]

Back-substitution:
\[
x_3 = \frac{-6}{-2} = 3, \quad
-2x_2 = -4 \implies x_2 = 2, \quad
x_1 + 2 + 3 = 6 \implies x_1 = 1.
\]

\textbf{Solution:} $x_1 = 1,\; x_2 = 2,\; x_3 = 3$.

\textit{Verification:}
\begin{align*}
\text{Eq.\,1:} & \quad 1 + 2 + 3 = 6. \checkmark \\
\text{Eq.\,2:} & \quad 1 - 2 + 3 = 2. \checkmark \\
\text{Eq.\,3:} & \quad 1 + 2 - 3 = 0. \checkmark
\end{align*}

All three original equations are satisfied.

%------------------------------------------------------------
\subsection{Gauss-Jordan Method}
%------------------------------------------------------------

Gauss elimination stops once the matrix is upper-triangular and then requires back-substitution.
Gauss-Jordan goes further: it eliminates \textit{above} each pivot as well, driving the matrix all the
way to Reduced Row Echelon Form (RREF).
In RREF, each pivot column has a single $1$ and zeros everywhere else --- the solution can be read
directly from the augmented column with no arithmetic required.

\begin{infobox}[Gauss-Jordan Algorithm and RREF]
\textbf{Reduced Row Echelon Form (RREF):} A matrix is in RREF when:
\begin{itemize}[leftmargin=1.5em]
  \item Each non-zero row has a leading $1$ (pivot).
  \item Every pivot column is a standard basis vector (zero except for the pivot entry).
  \item Pivot columns appear left-to-right in successive rows.
  \item Zero rows are at the bottom.
\end{itemize}

\textbf{Algorithm:} After forward elimination (same as Gauss), continue upward:
\begin{enumerate}[leftmargin=2em]
  \item \textit{Normalise:} Divide each pivot row by its pivot so the pivot becomes $1$.
  \item \textit{Upward elimination:} For each pivot in row $k$, eliminate all entries
        \textit{above} it using $R_i \leftarrow R_i - a_{ik} R_k$, $i = 1, \ldots, k-1$.
\end{enumerate}

\textbf{Key difference:}
\begin{center}
\begin{tabular}{@{}ll@{}}
\toprule
\textbf{Gauss} & Stops at upper triangular; requires back-substitution. \\
\textbf{Gauss-Jordan} & Continues to RREF; solution read directly from last column. \\
\bottomrule
\end{tabular}
\end{center}
\end{infobox}

\textbf{Worked Example A --- Gauss-Jordan (Unique Solution)}

Solve: $x + y + z = 6$, $x - y + z = 2$, $x + y - z = 0$ (same system as above).

\textit{After forward elimination} (from the Gauss example above):
\[
\left[\begin{array}{ccc|c}
1 & 1 & 1 & 6 \\
0 & -2 & 0 & -4 \\
0 & 0 & -2 & -6
\end{array}\right].
\]

\textit{Normalise rows:}

$R_2 \leftarrow R_2 / (-2)$: $[0, 1, 0, 2]$.\\
$R_3 \leftarrow R_3 / (-2)$: $[0, 0, 1, 3]$.

\[
\left[\begin{array}{ccc|c}
1 & 1 & 1 & 6 \\
0 & 1 & 0 & 2 \\
0 & 0 & 1 & 3
\end{array}\right].
\]

\textit{Upward elimination:}

Eliminate column 3 above $R_3$ (i.e., from $R_1$ and $R_2$):

$R_1 \leftarrow R_1 - 1 \cdot R_3$: $[1, 1, 0, 3]$.\\
$R_2$ has $0$ in column 3 already: no change.

\[
\left[\begin{array}{ccc|c}
1 & 1 & 0 & 3 \\
0 & 1 & 0 & 2 \\
0 & 0 & 1 & 3
\end{array}\right].
\]

Eliminate column 2 above $R_2$ (from $R_1$):

$R_1 \leftarrow R_1 - 1 \cdot R_2$: $[1, 0, 0, 1]$.

\textbf{RREF:}
\[
\left[\begin{array}{ccc|c}
1 & 0 & 0 & 1 \\
0 & 1 & 0 & 2 \\
0 & 0 & 1 & 3
\end{array}\right].
\]

\textbf{Solution read directly:} $x_1 = 1,\; x_2 = 2,\; x_3 = 3$. No back-substitution needed.

\begin{learnbox}[Gauss-Jordan Insight]
Gauss-Jordan gives the solution directly from the last column of the RREF --- no back-substitution needed.
\end{learnbox}

\textbf{Worked Example B --- Gauss-Jordan (Infinite Solutions)}

Solve the $3 \times 4$ system:
\[
x_1 + 2x_2 + x_3 + x_4 = 4, \quad
2x_1 + 4x_2 + 3x_3 + 2x_4 = 9, \quad
x_1 + 2x_2 + 2x_3 + 3x_4 = 7.
\]

Augmented matrix:
\[
\left[\begin{array}{cccc|c}
1 & 2 & 1 & 1 & 4 \\
2 & 4 & 3 & 2 & 9 \\
1 & 2 & 2 & 3 & 7
\end{array}\right].
\]

$m_{21}=2$: $R_2 \leftarrow R_2 - 2R_1$: $[0,0,1,0,1]$.\\
$m_{31}=1$: $R_3 \leftarrow R_3 - R_1$: $[0,0,1,2,3]$.

\[
\left[\begin{array}{cccc|c}
1 & 2 & 1 & 1 & 4 \\
0 & 0 & 1 & 0 & 1 \\
0 & 0 & 1 & 2 & 3
\end{array}\right].
\]

$m_{32}=1$: $R_3 \leftarrow R_3 - R_2$: $[0,0,0,2,2]$.\\
$R_3 \leftarrow R_3/2$: $[0,0,0,1,1]$.

Upward elimination for pivot in column 4:

$R_2$ has $0$ in col 4: no change.\\
$R_1 \leftarrow R_1 - 1\cdot R_3$: $[1,2,1,0,3]$.

Now pivot in column 3 (row 2):

$R_1 \leftarrow R_1 - 1\cdot R_2$: $[1,2,0,0,2]$.

\textbf{RREF:}
\[
\left[\begin{array}{cccc|c}
1 & 2 & 0 & 0 & 2 \\
0 & 0 & 1 & 0 & 1 \\
0 & 0 & 0 & 1 & 1
\end{array}\right].
\]

Pivot variables: $x_1, x_3, x_4$. Free variable: $x_2 = t \in \mathbb{R}$.

\begin{align*}
x_1 &= 2 - 2t, \\
x_2 &= t, \\
x_3 &= 1, \\
x_4 &= 1.
\end{align*}

\textbf{General solution in parametric vector form:}
\[
\mathbf{x} =
\begin{pmatrix} 2 \\ 0 \\ 1 \\ 1 \end{pmatrix}
+ t
\begin{pmatrix} -2 \\ 1 \\ 0 \\ 0 \end{pmatrix},
\qquad t \in \mathbb{R}.
\]

This represents a line through $(2,0,1,1)^T$ in the direction $(-2,1,0,0)^T$ in $\mathbb{R}^4$.

\begin{learnbox}[Infinite Solutions]
Infinite solutions mean free variables exist --- always express the complete general solution as
\(\mathbf{x} = \mathbf{x}_{\text{particular}} + t\,\mathbf{x}_{\text{homogeneous}}\).
\end{learnbox}

%============================================================
\section{Visual Artifacts}
%============================================================

\subsection*{Visual 1: Geometric Interpretation --- Three Cases of Two-Line Systems}

\begin{tikzpicture}
\begin{axis}[
  name=ax1,
  width=5cm, height=5cm,
  xmin=-1, xmax=4, ymin=-1, ymax=4,
  xlabel={$x$}, ylabel={$y$},
  title={Case 1: Unique (Intersect)},
  grid=major, grid style={dashed,gray!40},
  axis lines=center
]
\addplot[thick, blue, domain=-1:4] {x+1} node[right,font=\tiny]{$L_1$};
\addplot[thick, red,  domain=-1:4] {-x+3} node[right,font=\tiny]{$L_2$};
\addplot[only marks, mark=*, mark size=2pt, black] coordinates {(1,2)};
\node[font=\tiny,above right] at (axis cs:1,2) {$(1,2)$};
\end{axis}

\begin{axis}[
  name=ax2,
  at={($(ax1.east)+(1.5cm,0)$)}, anchor=west,
  width=5cm, height=5cm,
  xmin=-1, xmax=4, ymin=-1, ymax=4,
  xlabel={$x$}, ylabel={$y$},
  title={Case 2: Inconsistent (Parallel)},
  grid=major, grid style={dashed,gray!40},
  axis lines=center
]
\addplot[thick, blue, domain=-1:4] {x+2} node[right,font=\tiny]{$L_1$};
\addplot[thick, red,  domain=-1:4] {x-1} node[right,font=\tiny]{$L_2$};
\node[font=\tiny] at (axis cs:1.5,0) {No intersection};
\end{axis}

\begin{axis}[
  name=ax3,
  at={($(ax2.east)+(1.5cm,0)$)}, anchor=west,
  width=5cm, height=5cm,
  xmin=-1, xmax=4, ymin=-1, ymax=4,
  xlabel={$x$}, ylabel={$y$},
  title={Case 3: Infinite (Same Line)},
  grid=major, grid style={dashed,gray!40},
  axis lines=center
]
\addplot[ultra thick, blue!50, domain=-1:4] {x+1} node[right,font=\tiny]{$L_1 \equiv L_2$};
\addplot[thick, dashed, red, domain=-1:4] {x+1};
\node[font=\tiny] at (axis cs:1.5,0.2) {Overlap};
\end{axis}
\end{tikzpicture}

\subsection*{Visual 2: Augmented Matrix Transformation Stages}

\begin{tikzpicture}[
  node distance=1.5cm,
  box/.style={draw, rounded corners, fill=blue!8, minimum width=3.2cm, minimum height=2.6cm,
              align=center, font=\small},
  arr/.style={->, thick, >=stealth}
]

\node[box] (A) {$\left[\begin{smallmatrix} \mathbf{1} & 1 & 1 & 6 \\ 1 & -1 & 1 & 2 \\ 1 & 1 & -1 & 0 \end{smallmatrix}\right]$\\[4pt] Original $[\mathbf{A}|\mathbf{b}]$};

\node[box, right=of A] (B) {$\left[\begin{smallmatrix} \mathbf{1} & 1 & 1 & 6 \\ 0 & \mathbf{-2} & 0 & -4 \\ 0 & 0 & -2 & -6 \end{smallmatrix}\right]$\\[4pt] After Fwd.\ Elim.};

\node[box, right=of B] (C) {$\left[\begin{smallmatrix} \mathbf{1} & 0 & 0 & 1 \\ 0 & \mathbf{1} & 0 & 2 \\ 0 & 0 & \mathbf{1} & 3 \end{smallmatrix}\right]$\\[4pt] RREF};

\draw[arr] (A) -- node[above,font=\tiny]{Forward elim.} (B);
\draw[arr] (B) -- node[above,font=\tiny]{Backward elim. + normalise} (C);

\node[below=0.4cm of C, font=\small\itshape] {Solution: $(1,2,3)$};
\end{tikzpicture}

\subsection*{Visual 3: Rouch\'{e}--Capelli Decision Flowchart}

\begin{tikzpicture}[
  node distance=1.2cm and 1.8cm,
  decision/.style={diamond, draw, fill=orange!20, aspect=2, align=center, font=\small},
  process/.style={rectangle, draw, fill=blue!10, rounded corners, align=center, font=\small,
                  minimum width=3cm, minimum height=0.8cm},
  outcome/.style={rectangle, draw, fill=green!15, rounded corners, align=center, font=\small,
                  minimum width=3.5cm, minimum height=1cm},
  arr/.style={->, thick, >=stealth}
]

\node[process] (start) {Given $\mathbf{A}\mathbf{x}=\mathbf{b}$};
\node[process, below=of start] (ranks) {Compute $\rho_A = \text{rank}(\mathbf{A})$\\$\rho_{Ab} = \text{rank}([\mathbf{A}|\mathbf{b}])$};
\node[decision, below=of ranks] (eq) {$\rho_A = \rho_{Ab}$?};
\node[outcome, below left=of eq, xshift=-1cm] (incon) {\textbf{Inconsistent}\\No solution\\(Rouch\'{e}--Capelli)};
\node[decision, below right=of eq, xshift=1cm] (neqn) {$\rho_A = n$?};
\node[outcome, below left=of neqn] (uniq) {\textbf{Unique Solution}\\$\rho = n$};
\node[outcome, below right=of neqn] (inf)  {\textbf{Infinite Solutions}\\$(n-\rho)$ free vars};

\draw[arr] (start) -- (ranks);
\draw[arr] (ranks) -- (eq);
\draw[arr] (eq) -- node[left,font=\tiny]{No} (incon);
\draw[arr] (eq) -- node[right,font=\tiny]{Yes} (neqn);
\draw[arr] (neqn) -- node[left,font=\tiny]{Yes} (uniq);
\draw[arr] (neqn) -- node[right,font=\tiny]{No} (inf);
\end{tikzpicture}

\subsection*{Visual 4: Gauss vs.\ Gauss-Jordan --- Side by Side Reduction}

\begin{tikzpicture}[
  node distance=0.9cm,
  mat/.style={draw, fill=gray!8, rounded corners, align=center, font=\small,
              minimum width=3.6cm, minimum height=1.4cm},
  arr/.style={->, thick, >=stealth},
  label/.style={font=\small\bfseries}
]

\node[label] (g) {GAUSS};
\node[mat, below=0.4cm of g] (g0)
  {$\left[\begin{smallmatrix}1&1&1&6\\1&-1&1&2\\1&1&-1&0\end{smallmatrix}\right]$};
\node[mat, below=0.3cm of g0] (g1)
  {$\left[\begin{smallmatrix}1&1&1&6\\0&-2&0&-4\\0&0&-2&-6\end{smallmatrix}\right]$\\{\tiny Upper triangular --- STOP}};
\node[below=0.2cm of g1, font=\tiny] {Back-substitute to get $(1,2,3)$};

\node[label, right=4cm of g] (gj) {GAUSS-JORDAN};
\node[mat, below=0.4cm of gj] (gj0)
  {$\left[\begin{smallmatrix}1&1&1&6\\1&-1&1&2\\1&1&-1&0\end{smallmatrix}\right]$};
\node[mat, below=0.3cm of gj0] (gj1)
  {$\left[\begin{smallmatrix}1&1&1&6\\0&-2&0&-4\\0&0&-2&-6\end{smallmatrix}\right]$\\{\tiny Upper triangular --- CONTINUE}};
\node[mat, below=0.3cm of gj1] (gj2)
  {$\left[\begin{smallmatrix}1&0&0&1\\0&1&0&2\\0&0&1&3\end{smallmatrix}\right]$\\{\tiny RREF --- READ OFF}};
\node[below=0.2cm of gj2, font=\tiny] {Solution directly: $(1,2,3)$};

\draw[arr] (g0) -- (g1);
\draw[arr] (gj0) -- (gj1);
\draw[arr] (gj1) -- (gj2);
\end{tikzpicture}

%============================================================
\section{Step-by-Step Algorithmic Workflows}
%============================================================

\subsection*{Algorithm 1: Consistency Check via Rouch\'{e}--Capelli}

\begin{enumerate}[leftmargin=2em,label=\textbf{Step \arabic*.}]
  \item Write the coefficient matrix $\mathbf{A}$ and the augmented matrix $[\mathbf{A}|\mathbf{b}]$.
  \item Row-reduce $[\mathbf{A}|\mathbf{b}]$ to echelon form.
  \item Count the number of non-zero rows in $\mathbf{A}$: this is $\rho_A = \text{rank}(\mathbf{A})$.
  \item Count the number of non-zero rows in $[\mathbf{A}|\mathbf{b}]$: this is $\rho_{Ab} = \text{rank}([\mathbf{A}|\mathbf{b}])$.
  \item Apply the \textbf{Rouch\'{e}--Capelli theorem}:
    \begin{itemize}[leftmargin=1.5em]
      \item $\rho_A \neq \rho_{Ab}$: \textbf{Inconsistent.} Stop.
      \item $\rho_A = \rho_{Ab} = n$: \textbf{Unique solution.} Proceed to solve.
      \item $\rho_A = \rho_{Ab} = r < n$: \textbf{Infinite solutions.} Identify $(n-r)$ free variables.
    \end{itemize}
\end{enumerate}

\subsection*{Algorithm 2: Gauss Elimination}

\begin{enumerate}[leftmargin=2em,label=\textbf{Step \arabic*.}]
  \item Form augmented matrix $[\mathbf{A}|\mathbf{b}]$.
  \item \textit{Forward sweep:} For $k = 1$ to $n-1$:
    \begin{enumerate}[label=(\alph*)]
      \item If $a_{kk} = 0$, swap with a lower row (partial pivoting).
      \item For $i = k+1$ to $n$: compute $m_{ik} = a_{ik}/a_{kk}$ and apply
            $R_i \leftarrow R_i - m_{ik} R_k$.
    \end{enumerate}
  \item \textit{Back-substitution:} Solve from $x_n$ downward:
        $x_i = \bigl(b_i - \sum_{j=i+1}^n a_{ij} x_j\bigr) / a_{ii}$.
\end{enumerate}

\subsection*{Algorithm 3: Gauss-Jordan Elimination}

\begin{enumerate}[leftmargin=2em,label=\textbf{Step \arabic*.}]
  \item Perform Gauss forward elimination to reach upper-triangular form.
  \item \textit{Normalise:} For $k = n$ down to $1$, divide $R_k$ by $a_{kk}$ so the pivot becomes $1$.
  \item \textit{Backward sweep:} For each pivot row $k$, eliminate entries \textit{above}:
        For $i = 1$ to $k-1$: $R_i \leftarrow R_i - a_{ik} R_k$.
  \item Read off the solution directly from the last column of the RREF.
\end{enumerate}

\subsection*{Algorithm 4: General Solution for Infinite Solutions}

\begin{enumerate}[leftmargin=2em,label=\textbf{Step \arabic*.}]
  \item Identify pivot columns (basic variables) and free columns (free variables).
  \item Assign parameters $t_1, t_2, \ldots, t_{n-r}$ to the free variables.
  \item Express each basic variable in terms of the free variables.
  \item Write the complete solution as:
        $\mathbf{x} = \mathbf{x}_p + t_1 \mathbf{v}_1 + t_2 \mathbf{v}_2 + \cdots + t_{n-r} \mathbf{v}_{n-r}$.
\end{enumerate}

%============================================================
\section{Fully Worked Numerical Examples}
%============================================================

\subsection*{Example 1 --- Three Consistency Cases via Rouch\'{e}--Capelli}

(Presented fully in Section 3.1 above, with all three cases A, B, C.
We restate conclusions here for clarity.)

\begin{itemize}[leftmargin=1.5em]
  \item \textbf{Case A} ($2\times 2$, unique): $\rho_A = \rho_{Ab} = 2 = n$ $\Rightarrow$ unique solution.
  \item \textbf{Case B} ($2\times 2$, infinite): $\rho_A = \rho_{Ab} = 1 < 2$ $\Rightarrow$ infinite solutions, 1 free variable.
  \item \textbf{Case C} ($2\times 2$, inconsistent): $\rho_A = 1 \neq 2 = \rho_{Ab}$ $\Rightarrow$ inconsistent.
\end{itemize}

\begin{learnbox}[Rouch\'{e}--Capelli First]
Always apply the Rouch\'{e}--Capelli theorem first --- it tells you immediately whether solving is
worthwhile, saving time and avoiding meaningless computation.
\end{learnbox}

\subsection*{Example 2 --- Gauss Elimination with Unique Solution and Verification}

Solve $x + y + z = 6$, $x - y + z = 2$, $x + y - z = 0$.
(Full solution with every arithmetic step and verification of all three equations
is shown in Section 3.2.)

\textbf{Summary:} $x_1 = 1,\; x_2 = 2,\; x_3 = 3$, verified in all three original equations. $\checkmark$

\begin{learnbox}[Back-Substitution Order]
Back-substitution starts from the \textsc{last} equation and works upward, computing
$x_n$ first, then $x_{n-1}$, down to $x_1$.
\end{learnbox}

\subsection*{Example 3 --- Gauss-Jordan with Unique Solution}

(Full Gauss-Jordan reduction of the same system shown in Section 3.3, Worked Example A.)

\textbf{Key observation:} From forward-elimination onwards, Gauss-Jordan goes further by normalising
pivots and eliminating above each pivot.
The final RREF
\[
\left[\begin{array}{ccc|c} 1&0&0&1 \\ 0&1&0&2 \\ 0&0&1&3 \end{array}\right]
\]
gives $x_1=1, x_2=2, x_3=3$ with no additional arithmetic.

\begin{learnbox}[RREF Direct Read-Off]
Gauss-Jordan gives the solution directly from the last column of the RREF ---
no back-substitution needed.
\end{learnbox}

\subsection*{Example 4 --- Gauss-Jordan with Infinite Solutions and Power Network Application}

(Full solution with parametric vector form shown in Section 3.3, Worked Example B.)

\textbf{Engineering context:} In a 4-bus power distribution network with 3 KCL equations,
the bus voltages satisfy a $3 \times 4$ admittance system.
Finding that $x_2$ is a free variable corresponds to one bus voltage being uncontrolled
(a floating bus), meaning additional constraints (e.g., a voltage regulator) must be imposed.

\textbf{General solution:}
$\mathbf{x} = (2,0,1,1)^T + t(-2,1,0,0)^T,\; t \in \mathbb{R}$.

\begin{learnbox}[Parametric Vector Form]
Infinite solutions mean free variables exist --- always express the complete general solution as
$\mathbf{x} = \mathbf{x}_{\text{particular}} + t\,\mathbf{x}_{\text{homogeneous}}$.
\end{learnbox}

%============================================================
\section{Tabular Reference}
%============================================================

\begin{center}
\textbf{Table 1: Gauss vs.\ Gauss-Jordan Method Comparison}
\vspace{0.4em}

\begin{tabular}{@{}p{2.5cm}p{3cm}p{2.5cm}p{2.8cm}p{3cm}@{}}
\toprule
\textbf{Method} & \textbf{Reduction Target} & \textbf{Back-Sub Needed?} & \textbf{Computational Cost} & \textbf{Best For} \\
\midrule
Gauss Elimination & Upper triangular form & Yes & $\approx n^3/3$ flops & Single RHS, large sparse systems \\
\addlinespace
Gauss-Jordan & Reduced Row Echelon Form (RREF) & No & $\approx n^3/2$ flops & Matrix inverse, null space, multiple RHS \\
\bottomrule
\end{tabular}
\end{center}

\vspace{1em}

\begin{center}
\textbf{Table 2: Rouch\'{e}--Capelli Theorem --- All Three Cases}
\vspace{0.4em}

\begin{tabular}{@{}p{3cm}p{2.5cm}p{2.5cm}p{2.8cm}p{3cm}@{}}
\toprule
\textbf{Rank Condition} & \textbf{Solution Type} & \textbf{Free Variables} & \textbf{Geometric Meaning} & \textbf{Engineering Example} \\
\midrule
$\rho_A = \rho_{Ab} = n$ & Unique & 0 & Lines/planes intersect at one point & Single displacement field in FEA \\
\addlinespace
$\rho_A = \rho_{Ab} = r < n$ & Infinite & $n - r$ & Equations describe a subspace & Underdetermined circuit with floating nodes \\
\addlinespace
$\rho_A \neq \rho_{Ab}$ & None (Inconsistent) & N/A & Lines/planes are parallel without intersection & Contradictory boundary conditions in BVP \\
\bottomrule
\end{tabular}
\end{center}

%============================================================
\section{Common Student Mistakes and Pitfalls}
%============================================================

\begin{mistakebox}[Common Mistakes in Systems of Linear Equations]
\begin{tabular}{@{}p{4cm}p{5cm}p{5cm}@{}}
\toprule
\textbf{Mistake} & \textbf{What Goes Wrong} & \textbf{Correct Approach} \\
\midrule
Skipping consistency check & Solving a system that has no solution wastes all effort. &
  Apply the Rouch\'{e}--Capelli theorem \textit{first}: compute both ranks, compare. \\
\addlinespace
Confusing zero row with inconsistency & A row of all zeros in $[\mathbf{A}|\mathbf{b}]$ is fine
  (redundant equation). &
  Inconsistency only occurs when a row looks like $[0\ 0\ \cdots\ 0\ |\ c]$ with $c \neq 0$,
  i.e., $0 = c$. \\
\addlinespace
Arithmetic error in multiplier & Writing $m_{ik}$ incorrectly flips signs and corrupts all subsequent rows. &
  Write $m_{ik} = a_{ik}/a_{kk}$ explicitly before applying $R_i \leftarrow R_i - m_{ik}R_k$. \\
\addlinespace
Wrong back-substitution order & Starting from $x_1$ instead of $x_n$ gives circular dependencies. &
  Always start from the \textit{last} equation: solve $x_n$, then $x_{n-1}$, \ldots, then $x_1$. \\
\addlinespace
Calling upper-triangular RREF & Stopping at upper-triangular form and claiming Gauss-Jordan is done. &
  RREF requires both: (a) pivot $= 1$, and (b) all entries above and below each pivot $= 0$. \\
\addlinespace
Forgetting to verify & Arithmetic errors go undetected; wrong answers submitted with confidence. &
  Always substitute the final solution back into \textit{every} original equation and check. \\
\bottomrule
\end{tabular}
\end{mistakebox}

%============================================================
\section{Comprehensive Assessment Suite}
%============================================================

\subsection*{Viva-Voce Questions}

\begin{enumerate}[leftmargin=2em, label=\textbf{V\arabic*.}]
  \item State the Rouch\'{e}--Capelli theorem by name and write its formal condition for consistency.
  \item What does $\text{rank}(\mathbf{A}) \neq \text{rank}([\mathbf{A}|\mathbf{b}])$ geometrically mean
        for a system of two equations in two unknowns?
  \item If a $4 \times 6$ system has $\text{rank}(\mathbf{A}) = \text{rank}([\mathbf{A}|\mathbf{b}]) = 3$,
        how many free variables does the solution set have?
  \item Describe all the steps of Gauss elimination from forming the augmented matrix to writing the final solution.
  \item Why is partial pivoting used in Gauss elimination? Give an example of what goes wrong without it.
  \item After computing the solution by Gauss elimination, how do you verify it? How many equations must be checked?
  \item What is the key difference between the stopping point of Gauss elimination and Gauss-Jordan elimination?
  \item Define Reduced Row Echelon Form. State all four conditions a matrix must satisfy to be in RREF.
  \item A system has infinitely many solutions with $r = 2$ and $n = 5$. How do you express the general solution in parametric vector form?
\end{enumerate}

\subsection*{Descriptive Problems}

\begin{enumerate}[leftmargin=2em, label=\textbf{D\arabic*.}]
  \item Apply the Rouch\'{e}--Capelli theorem to each of the following systems and state the nature of the
        solution without solving:
        (a) $x + y = 3$, $2x + 2y = 6$;\quad
        (b) $x + y = 3$, $2x + 2y = 7$;\quad
        (c) $x + y + z = 3$, $2x - y + z = 0$, $x + 2y + 3z = 5$.

  \item Solve the system $2x + y - z = 8$, $-3x - y + 2z = -11$, $-2x + y + 2z = -3$ using
        Gauss elimination. Show all multipliers explicitly, perform back-substitution from the
        last equation upward, and verify your solution in all three original equations.

  \item Using the Gauss-Jordan method, find the complete solution to
        $x_1 + 2x_2 - x_3 = 3$, $2x_1 + 5x_2 + 2x_3 = 0$, $x_1 + 4x_2 + 7x_3 = -9$.
        Display each augmented matrix after every elementary row operation and read the solution
        from the final RREF.

  \item A thermal network yields the $3 \times 5$ system:
        $T_1 + T_2 + T_3 + T_4 + T_5 = 100$,
        $2T_1 + T_2 + 3T_3 - T_4 + T_5 = 150$,
        $T_1 + 2T_2 + T_3 + 2T_4 - T_5 = 80$.
        Apply the Rouch\'{e}--Capelli theorem, then use Gauss-Jordan to find the general solution
        in parametric vector form. Interpret the free variables in the context of unspecified node temperatures.
\end{enumerate}

\subsection*{Multiple Choice Questions}

\begin{enumerate}[leftmargin=2em, label=\textbf{MCQ\arabic*.}]
  \item A system $\mathbf{A}\mathbf{x} = \mathbf{b}$ with $m = 4$ equations and $n = 4$ unknowns has
        $\text{rank}(\mathbf{A}) = 3$ and $\text{rank}([\mathbf{A}|\mathbf{b}]) = 4$. The system is:
        \begin{enumerate}[label=(\alph*)]
          \item Consistent with a unique solution
          \item Consistent with infinitely many solutions
          \item \textbf{Inconsistent (no solution)}
          \item Consistent with exactly 4 solutions
        \end{enumerate}
        \textit{Explanation:} $\rho_A \neq \rho_{Ab}$, so by the Rouch\'{e}--Capelli theorem the system is inconsistent.

  \item For a consistent $5 \times 8$ system with $\text{rank}(\mathbf{A}) = 4$, the number of free variables is:
        \begin{enumerate}[label=(\alph*)]
          \item 1
          \item 3
          \item \textbf{4}
          \item 5
        \end{enumerate}
        \textit{Explanation:} Free variables $= n - r = 8 - 4 = 4$.

  \item Gauss elimination transforms the augmented matrix into:
        \begin{enumerate}[label=(\alph*)]
          \item Reduced Row Echelon Form
          \item \textbf{Upper triangular form}
          \item Diagonal form
          \item Identity form
        \end{enumerate}
        \textit{Explanation:} Gauss stops at upper triangular; Gauss-Jordan continues to RREF.

  \item In Gauss elimination the multiplier for eliminating entry $a_{32}$ using pivot $a_{22}$ is:
        \begin{enumerate}[label=(\alph*)]
          \item $a_{22}/a_{32}$
          \item \textbf{$a_{32}/a_{22}$}
          \item $a_{32} - a_{22}$
          \item $a_{22} - a_{32}$
        \end{enumerate}
        \textit{Explanation:} $m_{32} = a_{32}/a_{22}$; the operation is $R_3 \leftarrow R_3 - m_{32}R_2$.

  \item Which statement correctly distinguishes Gauss-Jordan from Gauss elimination?
        \begin{enumerate}[label=(\alph*)]
          \item Gauss-Jordan uses fewer arithmetic operations.
          \item Gauss-Jordan requires back-substitution; Gauss does not.
          \item \textbf{Gauss-Jordan eliminates above and below each pivot; Gauss only eliminates below.}
          \item Both methods produce the same intermediate matrix at every step.
        \end{enumerate}
        \textit{Explanation:} Gauss-Jordan performs full reduction (above \textit{and} below) to reach RREF,
        so no back-substitution is needed.

  \item A row of the form $[0 \ 0 \ 0 \ | \ 5]$ appearing in the row-reduced augmented matrix indicates:
        \begin{enumerate}[label=(\alph*)]
          \item A redundant equation (consistent)
          \item A free variable
          \item \textbf{The system is inconsistent}
          \item The system has infinitely many solutions
        \end{enumerate}
        \textit{Explanation:} This row represents $0 = 5$, a contradiction, making the system inconsistent.
\end{enumerate}

%============================================================
\section{Quick Recap and Formula Sheet}
%============================================================

\begin{learnbox}[Key Formulas and Principles --- System of Linear Equations]
\begin{itemize}[leftmargin=1.5em]
  \item \textbf{Rouch\'{e}--Capelli Theorem:} $\mathbf{A}\mathbf{x} = \mathbf{b}$ is consistent
        $\Leftrightarrow$ $\text{rank}(\mathbf{A}) = \text{rank}([\mathbf{A}|\mathbf{b}])$.
  \item \textbf{Unique solution:} $\text{rank}(\mathbf{A}) = \text{rank}([\mathbf{A}|\mathbf{b}]) = n$.
  \item \textbf{Infinite solutions:} $\text{rank}(\mathbf{A}) = \text{rank}([\mathbf{A}|\mathbf{b}]) = r < n$
        $\Rightarrow$ $(n-r)$ free variables.
  \item \textbf{Inconsistent:} $\text{rank}(\mathbf{A}) \neq \text{rank}([\mathbf{A}|\mathbf{b}])$
        (row $[0\ \cdots\ 0\ |\ c]$, $c \neq 0$ appears in reduction).
  \item \textbf{Gauss elimination:} Forward sweep using $m_{ik} = a_{ik}/a_{kk}$,
        operation $R_i \leftarrow R_i - m_{ik}R_k$; then back-substitute from last equation upward.
  \item \textbf{Gauss-Jordan:} Forward + backward sweeps $\Rightarrow$ RREF; solution read directly
        from last column --- no back-substitution needed.
  \item \textbf{Row operation notation:} $R_i \leftarrow R_i - m_{ik}R_k$,\quad $m_{ik} = a_{ik}/a_{kk}$.
  \item \textbf{General solution (infinite case):}
        $\mathbf{x} = \mathbf{x}_{\text{particular}} + t_1\mathbf{v}_1 + \cdots + t_{n-r}\mathbf{v}_{n-r}$.
  \item \textbf{Always verify:} Substitute the final solution back into \textit{every} original equation
        and confirm equality --- never skip this step.
\end{itemize}
\end{learnbox}

\end{document}
